\chapter{引言}\label{ch:introduction}

\section{科研数据整理的真实困境}

生物医学研究中，“找到数据”不等于找到能用的数据。可以直接分析的数据集至少要说明：覆盖哪些实体和样本，来自哪个版本，原始文件如何转换为当前字段，不同来源能否合并，哪些值存在冲突，哪些判断需要研究者确认。只依靠对话和临时文件，难以逐项核对这些问题。

第一，同类数据也可能无法直接合并。相同列名未必代表相同测量；探针级与基因级表达表具有不同观测单位，原始强度、归一化值与对数值不能混用。生物活性数据还需对齐化合物标识、实验类型和单位，并记录冲突。

第二，可读信息不等于可计算数据。PDF 表格、扫描图、折线图与补充 XLSX 需要不同解析方法。视觉模型读取结果需附上页码、位置、模型名、提示版本与置信度，才能进一步复核。

第三，复制整理容易遗漏处理记录。在浏览器、Excel、脚本与聊天窗口之间反复搬运数据，可能丢失来源版本、筛选条件、字段映射与冲突处理依据，出现错误时难以回查。

第四，长任务需要保存运行状态。下载、权限审批和故障恢复可能跨越多个会话；仅依靠对话上下文，重启后容易重复操作或遗忘约束。

第五，来源分散且访问方式不同。同一问题可能涉及 PubMed、GEO、ClinVar、PDB 与论文附件。检索结果相关，不代表已经遍历分页、检查补充材料或覆盖全部必要来源。

\section{相关工作的实际缺口}\label{sec:related}

现有工作分别提供了检索、解析、数据整合和智能体编排能力。本文关注的是：这些环节产生的数据进入正式交付前，是否留下可逐项核对的来源、处理记录和发布检查结果。

数据库与解析工具解决查找和读取问题。PubMed、GEO、GDC、UniProt、ClinVar、ChEMBL 等提供专门查询入口；规则解析器适合固定格式，光学字符识别（Optical Character Recognition，OCR）与视觉语言模型（Vision-Language Model，VLM）用于读取页面和图像。跨库查询和字段对应仍需按任务设计。

数据整合与溯源工具提供了组织处理步骤和记录来源的方法。抽取—转换—加载（Extract–Transform–Load，ETL）流程适合执行明确的数据转换；W3C PROV、FAIR 和 OpenLineage 则提供来源描述或数据管理原则。应用仍需决定何时检查、如何拒绝不合格结果，以及怎样保留修正前后的文件。

LLM 编排与工作流框架可以支持工具调用、状态保存和人工介入，例如 LangGraph、AutoGen、OpenAI Agents SDK 与 Temporal\cite{langgraph,autogen,openai-agents,temporal-approval}。DocETL 提供查询改写与任务相关评估，Palimpzest 则在时长、费用和输出质量之间选择执行计划\cite{docetl,palimpzest}。这些能力可以用于构建发布检查，但框架支持某种检查，并不等于具体应用已经实现和验证了完整的数据交付流程。

本文将这些能力用于科研数据交付：来源登记后生成候选，候选经过检查与必要审批后正式发布，读取时再次核验。“信任边界”在本文中特指这条发布检查通道，不表示系统能自动证明科学内容正确。

\section{本系统的核心目标}

\mbox{BioMed QAgent} 聚焦于数据获取与处理，不替研究者作出科学结论。系统将自然语言需求转换为带有表结构、来源和检查记录的数据表，供后续分析使用，并支持根据反馈修正。

系统依次完成四项工作：

\begin{enumerate}
  \item 找得到：根据研究目标查找论文、数据库、附件、表格或图像。以 EGFR 突变体抑制实验为例，Agent 需要定位相关数据库和补充材料，并记录检索范围，而不只是返回相关网页。需求编译与来源发现见\ref{sec:requirement-compile} 节和\ref{sec:source-discovery} 节。
  \item 读得懂：将不同载体解析为字段含义明确的数据记录。例如表达表必须说明一行代表探针还是基因，不能仅凭列名推断。解析与规范化过程见第四章。
  \item 合得上：对齐实体标识和可换算的单位，同时保留缺失、重复和冲突。同一化合物的不同数据库标识需要映射，nM 与 $\mu$M 等单位差异需要按明确规则换算。
  \item 查得清：记录来源、输入摘要、处理版本、检查结果和人工修正。无法确认的内容应排除、标注或交研究者判断；正式发布与原始需求是否全部满足分别检查。
\end{enumerate}

上述流程对应来源发现、异构解析、清洗整合、图表提取与来源追溯。科研分析和结论生成不计入方向1A的正式交付要求。

\section{研究问题与数据需求}

本项目研究如何将自然语言数据需求转换为结构明确、处理步骤可检查的数据表。数据获取与整合属于本系统任务，后续假设验证与科学结论不属于正式交付范围。

用户输入自由文本后，系统从中确定数据集的十项选择，包括需要哪些数据、一行代表什么、从哪里采集、如何处理冲突等（表~\ref{tab:spec-choices}）。这些选择组成机器可读的数据集执行规格。规格确定了本次构建范围，但仍须与原始需求逐项核对，不能仅因规格执行成功就认为任务全部完成。

\begin{table}[htbp]
\centering
\caption{数据集规格的十项关键选择}\label{tab:spec-choices}
\begin{tabularx}{\linewidth}{>{\hsize=0.60000\hsize}L>{\hsize=1.20000\hsize}L>{\hsize=1.20000\hsize}L}
\toprule
\textbf{选择} & \textbf{含义} & \textbf{改变意味着什么} \\
\midrule
数据产品类别 & 本次需要哪一类数据产品 & 决定交付物的整体形态 \\
行层级 & 每行代表什么观测或统计单位 & 决定哪些行可合并、哪些需拆表 \\
目标表结构 & 交付表包含哪些字段及关系 & 决定字段映射、类型与检查规则 \\
实体集合 & 涉及哪些对象、范围多大 & 决定采集与筛选范围 \\
队列与筛选条件 & 按什么标准纳入或排除样本 & 改变样本总体与可比性 \\
来源与采集绑定 & 从哪个数据库或文件获取数据 & 决定可核对的来源 \\
规范化配置 & 采用哪些转换规则及参数 & 决定值的含义能否保持一致 \\
合并策略 & 重叠或冲突信息如何取舍 & 决定保留结果与冲突记录 \\
验证配置 & 采用哪些规则与阈值 & 决定发布检查标准 \\
输出格式 & 以什么格式交付 & 决定下游读取方式 \\
\bottomrule
\end{tabularx}
\end{table}

规格、来源文件、处理代码和参数均固定后，才有条件检查重复执行是否产生相同结果。重新发起来源检索或模型抽取，仍可能得到不同输入；样例1在 flash 与 max 下纳入研究不同，就是这种差异。第五章的哈希重验检查既有文件是否与清单一致，并非独立重跑实验。

相应地，系统交付的也不是一张孤立的数据表。合格的候选结果必须同时包含六类信息：数据本身、表结构、溯源信息、审计记录、验证结果、文件清单。候选结果只有通过全部质量检查后，系统才会生成正式发布版。正式发布只新增版本、不覆盖既有发布；读取时，系统依据发布清单重新校验文件摘要，发现文件缺失、损坏或与清单记录不一致即拒绝展示。

\section{主要贡献与核心机制}\label{sec:contributions}

本文提出三项相互配合的系统贡献，共同使用 Dataset Core 的八阶段流程。内容摘要、事件日志、模式检查和人工审批均为既有方法；本文将它们组合到同一次科研数据任务中，使候选文件在交付前接受统一检查。

\textbf{以可核验文件判定正式发布。}系统不将模型自报、工具预览或工作区 CSV 视为正式结果。Agent 提交规格，Core 完成处理与检查后生成发布清单。运行时向 Agent 提供真实进度，使其区分“仍在执行”“未能发布”和“已有正式发布物”。发布状态说明文件是否通过配置的检查；需求完成情况另按原始任务核对。

\textbf{检查运行时提出的表结构。}Agent 可以为未注册的多表任务提出表、字段、主外键与变换方案。候选文件必须说明输入来源，实际产生输出，并保证输出列与声明逐一对应。Core 重新计算文件摘要，再执行完整性检查、产品评估和发布检查；变换方不能自行选择或跳过这些模块。此处检查结构与来源是否一致，不自动判断变换计算是否正确。

\textbf{让人工确认对应具体证据。}审批请求携带任务身份、候选值与证据摘要。系统恢复操作前检查响应是否仍对应原请求和原对象；对象已变化或请求失效时，旧响应不能继续用于该对象。对已发起且要求人工处理的发布验收，只有人工接受后才继续发布。修正形成新版本，旧版文件与审批记录继续保留。各发布路径是否实际触发验收，须结合第五章运行记录判断。

来源文件、处理记录与发布清单共同说明数据从何而来、经过哪些步骤、最终交付哪些文件。摘要重验可以发现文件与记录不符，但摘要一致不等于数值正确，也不证明已经查全来源。这两项仍需对照原始来源和任务要求核验。

